\chapter{Swap Clearing and the Clearing-House Basis}\label{ch:m2:swap-clearing-and-the-ccp-basis}

In 2009 the leaders of the G20 agreed that every standardised over-the-counter
derivative should be cleared through a central counterparty by the end of
2012. Swaps that had been bilateral promises between banks became trades
with a clearing house, margined daily, and a swap's price was supposed to
depend only on its terms. It did not quite work out that way. By 2017 a
thirty-year dollar swap cleared at one clearing house was being quoted several
basis points away from the identical swap cleared at the other, and a data
firm could show that for USD~100 million the difference paid for about 727
thousand dollars of margin funding. The swap's cash flows are the same; what
differs is who else trades at each clearing house, and so how much margin a
dealer has to post to hold the other side. This chapter explains clearing,
the margin that comes with it, and why the same swap can have two prices.

\section{The clearing mandate}

\begin{definition}[Clearing mandate, client clearing and porting]\label{def:m2:swap-clearing-and-the-ccp-basis:mandate}
A \emph{clearing mandate}\index{clearing mandate} is a regulatory requirement
that designated classes of derivatives be cleared through a central
counterparty (One Quant Book 1, chapter 5). Firms that are not members of the
clearing house clear through a member, which guarantees their trades: this is
\emph{client clearing}\index{client clearing}. \emph{Porting}\index{porting} is
the transfer of a client's positions and margin from a defaulting clearing
member to another, so that the client's hedges survive its broker's failure.
\end{definition}

The United States implemented the commitment first: the Commodity Futures
Trading Commission adopted its first clearing requirement in November 2012, for
four classes of \omterm{def:m2:interest-rate-swaps:swap}{interest-rate swaps} and two of credit index swaps, and dealers
and active funds began clearing them in March 2013. The European Union followed
under its own regulation, and in 2024 added a requirement that its firms keep an
active account at a clearing house inside the Union, a political choice about
where euro swaps are cleared as much as a risk rule.

What clearing changes is the shape of the network
(\cref{fig:m2:swap-clearing-and-the-ccp-basis:network}). Five dealers trading
bilaterally face ten counterparty relationships, each with its own collateral
agreement and its own exposure; cleared, each faces one clearing house, its
trades with all the others net into one position there, and a default is
absorbed by the house's margin and default fund rather than passed along the
web.

\begin{omfigure}
\centering
\begin{tikzpicture}[font=\small,
  dl/.style={circle,draw=omDef,thick,fill=omDef!10,minimum size=7mm,inner sep=0pt}]
\foreach \i in {1,...,5} {\node[dl] (a\i) at ({90+72*(\i-1)}:1.6) {D\i};}
\foreach \i/\j in {1/2,1/3,1/4,1/5,2/3,2/4,2/5,3/4,3/5,4/5} {\draw[gray,thick] (a\i) -- (a\j);}
\node[font=\footnotesize,align=center] at (0,-2.5) {bilateral: 10 relationships};
\begin{scope}[xshift=6.5cm]
\node[rectangle,draw=omThm,thick,fill=omThm!12,minimum size=9mm,rounded corners=2pt] (c) at (0,0) {CCP};
\foreach \i in {1,...,5} {\node[dl] (b\i) at ({90+72*(\i-1)}:1.6) {D\i}; \draw[omThm,thick] (b\i) -- (c);}
\node[font=\footnotesize,align=center] at (0,-2.5) {cleared: 5 relationships};
\end{scope}
\end{tikzpicture}

\omcaption{Five dealers trading swaps bilaterally, each pair with its own
collateral agreement, and the same dealers after clearing, each facing only the
clearing house. Schematic.}\label{fig:m2:swap-clearing-and-the-ccp-basis:network}
\end{omfigure}

\begin{dated}{2026-09}{Mandates and margin rules}\label{dat:m2:swap-clearing-and-the-ccp-basis:mandates}
\textbf{United States}: clearing required for the designated swap classes from
11 March 2013 for dealers and active funds, later for others. \textbf{European
Union}: under EMIR 3, in force since 24 December 2024, firms subject to the
clearing obligation must hold an active account at an authorised EU clearing
house and clear a representative number of euro and zloty interest-rate
derivatives through it, from 24 June 2025. \textbf{Uncleared trades}: the
international margin framework reached its final phase on 1 September 2022, for
firms with more than EUR~8 billion of uncleared derivatives; initial margin need
not be exchanged below EUR~50 million per counterparty group.
\end{dated}

\section{The competing clearing houses}

Dollar swaps are cleared mainly at two houses, LCH's SwapClear and CME; other
currencies have their own mix, and part of euro clearing must now take place
inside the Union (\cref{dat:m2:swap-clearing-and-the-ccp-basis:mandates}). A dealer is a member of several and can choose where to clear
a trade with another dealer. A client usually clears where its clearing broker
and its own risk systems are set up, and it cannot move trades freely between
houses: margin is computed on each house's portfolio separately, and a trade
at one never offsets a trade at another.

\begin{omfigure}
\begin{tikzpicture}[font=\small,
  box/.style={draw,thick,rounded corners=2pt,align=center,minimum height=1cm,minimum width=2.9cm,inner sep=4pt},
  arr/.style={-{Stealth[length=2.5mm]},thick}]
\node[box,draw=omDef,fill=omDef!8] (cl) at (0,0) {clients\\(mostly pay fixed)};
\node[box,draw=omThm,fill=omThm!8] (dA) at (5,1.3) {dealer at CCP A\\receives fixed};
\node[box,draw=omThm,fill=omThm!8] (dB) at (5,-1.3) {dealer at CCP B\\pays fixed};
\node[box,draw=gray,fill=gray!8] (st) at (10.2,-1.3) {other dealers\\(receive fixed)};
\draw[arr] (cl) -- node[above left=2pt,align=center] {pay fixed\\at CCP A} (dA);
\draw[arr,dashed] (dA) -- node[right=4pt,align=left] {same dealer,\\hedges} (dB);
\draw[arr] (dB) -- node[below=2pt,align=center] {pays fixed\\at CCP B} (st);
\node[font=\footnotesize,text=omProp,align=center] at (10.2,1.3) {margin posted twice:\\one-way at A, one-way at B};
\end{tikzpicture}

\omcaption{How a clearing-house basis arises. Clients at one house are mostly
fixed payers; the dealers who take the other side receive fixed there and
hedge by paying fixed at the other house, where other dealers receive. The
dealer's risk nets to zero across the two, but each house margins its half on
its own, and the dealer funds both margins.}\label{fig:m2:swap-clearing-and-the-ccp-basis:flows}
\end{omfigure}

\section{Initial margin, cleared and uncleared}

A clearing house collects variation margin, the daily change in a swap's
value, and initial margin, a buffer against the loss it would suffer in closing
out a defaulted member's portfolio (One Quant Book 1, chapter 20).

\begin{definition}[Uncleared margin rules and credit support annex]\label{def:m2:swap-clearing-and-the-ccp-basis:umr}
The \emph{uncleared margin rules}\index{uncleared margin rules} are the
regulations requiring large users of derivatives that are not centrally
cleared to exchange initial and variation margin bilaterally, with the initial
margin held by a third party. A \emph{credit support annex}\index{credit support
annex} (CSA) is the part of a bilateral derivatives agreement that governs the
collateral: which assets, thresholds, minimum transfers, the rate paid on cash
collateral, and the timing of calls.
\end{definition}

\begin{proposition}[Initial margin and its cost]\label{prop:m2:swap-clearing-and-the-ccp-basis:mva}
Suppose a house margins a swap position at $z$ standard deviations of its value
change over a margin period of risk of $h$ days, with a daily rate volatility of
$\sigma$ basis points: $\mathrm{IM}(t) = z\,\sigma\sqrt{h}\,\mathrm{DV01}(t)$. If
posting margin costs a funding spread $s$ a year, the \emph{margin valuation
adjustment} is
\[
\mathrm{MVA} \;=\; \int_0^T s\,\mathrm{IM}(t)\,P(0,t)\,dt,
\]
and the running spread on the swap that pays for it is
$\mathrm{MVA}/\mathrm{DV01}(0)$ basis points.
\end{proposition}
\begin{proof}
The first statement defines the model: value changes are normal with
standard deviation $\sigma\sqrt h$ basis points of rate times the DV01. The cost
of holding $\mathrm{IM}(t)$ for $dt$ is $s\,\mathrm{IM}(t)\,dt$, discounted. A
running spread of $b$ basis points on the swap is worth $b\,\mathrm{DV01}(0)$
today.
\end{proof}

\begin{example}[A ten-year and a thirty-year swap]\label{ex:m2:swap-clearing-and-the-ccp-basis:mva}
With illustrative parameters, $\sigma = 7$ basis points a day, a five-day margin
period and $z = 2.326$, initial margin is 36.4 basis points of DV01. On USD~100
million of a ten-year swap at 4\%, the DV01 is USD~81\,109, initial margin starts at
USD~2.95 million and declines with the swap's remaining life
(\cref{fig:m2:swap-clearing-and-the-ccp-basis:im}); funded at 50 basis points a year
it costs USD~76\,442 over the swap's life: 0.94 basis points a year of running
spread. For the thirty-year swap, 2.30 basis points.
\end{example}

\begin{omfigure}
\begin{tikzpicture}
\begin{axis}[width=11cm,height=4.8cm,
  xlabel={years since the trade},ylabel={initial margin (USD million)},
  tick label style={font=\small},label style={font=\small},
  xmin=0,xmax=30,ymin=0,ymax=7,grid=major,grid style={gray!25},
  legend columns=2,legend style={font=\footnotesize,at={(0.5,-0.3)},anchor=north,draw=none,fill=none,/tikz/every even column/.append style={column sep=8pt}},
  table/col sep=comma]
\addplot[omDef,very thick,restrict x to domain=0:9.9] table[x=t,y=ten]{figdata/markets-2/10-swap-clearing-and-the-ccp-basis/im.csv};
\addplot[omThm,very thick] table[x=t,y=thirty]{figdata/markets-2/10-swap-clearing-and-the-ccp-basis/im.csv};
\legend{10-year swap,30-year swap}
\end{axis}
\end{tikzpicture}

\omcaption{Initial margin over the life of USD~100 million of a ten-year and a
thirty-year swap, in the stylised model of
\cref{prop:m2:swap-clearing-and-the-ccp-basis:mva}. Margin follows the DV01 of
the remaining swap, and the area under each curve, times the funding spread,
is its cost. Parameters are illustrative. Data: the chapter's
tutorial.}\label{fig:m2:swap-clearing-and-the-ccp-basis:im}
\end{omfigure}

\section{Why the same swap has two prices}

\begin{definition}[CCP basis]\label{def:m2:swap-clearing-and-the-ccp-basis:basis}
The \emph{CCP basis}\index{CCP basis} is the difference between the par rates
of economically identical swaps cleared at two different clearing houses.
\end{definition}

If a dealer's trades at one house net out, it posts little initial margin
there. If its clients at that house all pay fixed, it is a one-way receiver
there, and its hedges, paying fixed at the other house, leave it a one-way
payer there: two margins, both funded, for a position that is flat. The dealer
charges for this when it receives fixed at the first house, by asking a higher
fixed rate there. The analysis of 2017 found the effect largest in long dollar
swaps, where client payers at CME outnumbered receivers: about 3.4 basis points
at thirty years, of which 1.3 was the margin cost at CME and 2.1 at LCH.

\begin{omfigure}
\begin{tikzpicture}
\begin{axis}[width=11cm,height=4.8cm,
  xlabel={swap maturity (years)},ylabel={basis per one-way position (bp)},
  tick label style={font=\small},label style={font=\small},
  xmin=0,xmax=31,ymin=0,ymax=3.7,grid=major,grid style={gray!25},
  legend columns=3,legend style={font=\footnotesize,at={(0.5,-0.3)},anchor=north,draw=none,fill=none,/tikz/every even column/.append style={column sep=8pt}},
  table/col sep=comma]
\addplot[omDef,thick,dashed,mark=*,mark size=1.4pt] table[x=maturity,y=f25]{figdata/markets-2/10-swap-clearing-and-the-ccp-basis/basis.csv};
\addplot[omThm,very thick,mark=*,mark size=1.4pt] table[x=maturity,y=f50]{figdata/markets-2/10-swap-clearing-and-the-ccp-basis/basis.csv};
\addplot[omProp,thick,dotted,mark=*,mark size=1.4pt] table[x=maturity,y=f75]{figdata/markets-2/10-swap-clearing-and-the-ccp-basis/basis.csv};
\legend{funding 25 bp,funding 50 bp,funding 75 bp}
\end{axis}
\end{tikzpicture}

\omcaption{The running spread that pays for one one-way margined position, by
maturity and funding spread, in the stylised model. A dealer with one-way
positions at two houses needs twice as much. The basis grows with
maturity, because margin is held for longer, but less than in proportion,
because the margin runs off with the swap and later years are discounted. Data:
the chapter's tutorial.}\label{fig:m2:swap-clearing-and-the-ccp-basis:basis}
\end{omfigure}

The basis is therefore not an arbitrage. Receiving fixed at the rich house and
paying at the cheap one locks in the spread, but also creates the two
one-way positions whose margin cost the spread pays for; only a firm whose
existing positions at the two houses make the new trades reduce its margin
can earn it, and dealers compete to be that firm. Mandates that push flows to a
particular house, like the EU active account, move the basis too.

\section{Tutorial: pricing the margin}\label{tut:m2:swap-clearing-and-the-ccp-basis:mva}

\begin{tutorial}
\textbf{Goal.} Compute the initial-margin profile of a swap, its funding cost
and the basis that pays for it. \textbf{End state:}
\cref{fig:m2:swap-clearing-and-the-ccp-basis:im,fig:m2:swap-clearing-and-the-ccp-basis:basis}
and the numbers of \cref{ex:m2:swap-clearing-and-the-ccp-basis:mva}.
\begin{tutsteps}
\item \textbf{The margin model}: initial margin per unit of DV01.
\omcode{code/firm/ccpbasis/firm_ccpbasis.py}{24}{32}{A stylised initial-margin model.}\label{lst:m2:swap-clearing-and-the-ccp-basis:model}
\item \textbf{Margin over the life, its cost, and the basis.}
\omcode{code/firm/ccpbasis/firm_ccpbasis.py}{44}{58}{Initial-margin path, margin valuation adjustment, running basis.}\label{lst:m2:swap-clearing-and-the-ccp-basis:mva}
\item \textbf{Run} \texttt{ccp\_demo.one\_way\_book()} and
\texttt{fig\_ccp.py}.
\end{tutsteps}
\textbf{What to change next.} Double the margin period of risk to ten days, as
for an uncleared trade, and watch initial margin rise by $\sqrt 2$; then let the
funding spread fall to zero and explain why the basis disappears.
\end{tutorial}

\section{Build: the margin-funding calculator}\label{bld:m2:swap-clearing-and-the-ccp-basis:mva}

\begin{build}
\bfield{Purpose} The miniature firm quotes swaps to clients at different
clearing houses and must know what each trade costs it in margin at each
house; its risk system (One Quant Book 6) extends this to the real margin
models.
\bfield{Interface} \texttt{MarginModel(sigma\_bp\_day, mpor\_days, z)} with
\texttt{im\_per\_dv01()}; \texttt{annuity\_remaining(t, maturity, rate)};
\texttt{dv01\_path}; \texttt{im\_path}; \texttt{mva(notional, maturity, rate,
model, funding\_spread)}; \texttt{basis\_bp(\ldots, ccps)}.
\bfield{Rules} Flat curve, annual periods; margin proportional to the
remaining swap's DV01; quarterly steps for the funding integral; the basis per
one-way position, multiplied by the number of houses where the dealer is one
way.
\bfield{Acceptance tests} \texttt{code/firm/ccpbasis/tests/}: the annuity's
closed form; margin proportional to DV01 and declining; MVA linear in notional
and funding; basis increasing with maturity and additive across houses.
\bfield{Stretch} A portfolio margin (netting across trades at one house); a
historical-simulation margin on the curve of
\cref{bld:m2:interest-rate-swaps:curve}; the incremental margin of a new trade
at each house, which is what a dealer really prices.
\end{build}

\begin{omsources}
\item G20, \emph{Leaders' Statement: The Pittsburgh Summit}, September 2009.
\item Commodity Futures Trading Commission, press releases on the first
clearing determination (November 2012) and its phases (2013).
\item Basel Committee and IOSCO, \emph{Margin requirements for non-centrally
cleared derivatives}, and statements on the final phases (2019, 2020).
\item Clarus Financial Technology, ``CME-LCH basis for dummies'', June 2017; DLA
Piper, ``EMIR 3 active account requirement'', June 2025.
\end{omsources}

\section{Exercises}

\begin{exercise}[$\star$]\label{exo:m2:swap-clearing-and-the-ccp-basis:1}
In the model of \cref{ex:m2:swap-clearing-and-the-ccp-basis:mva}, what initial
margin does USD~100 million of a new ten-year swap require?
\end{exercise}

\begin{exercise}[$\star$]\label{exo:m2:swap-clearing-and-the-ccp-basis:2}
A client receives fixed on that swap. Rates rise 3 basis points in a day. What
variation margin does it pay?
\end{exercise}

\begin{exercise}[$\star$]\label{exo:m2:swap-clearing-and-the-ccp-basis:3}
A client's clearing broker defaults. What happens to the client's swaps if
\omterm{def:m2:swap-clearing-and-the-ccp-basis:mandate}{porting} succeeds, and if it fails?
\end{exercise}

\begin{exercise}[$\star\star$]\label{exo:m2:swap-clearing-and-the-ccp-basis:4}
Give the MVA of USD~100 million of the ten-year swap at a 50-basis-point funding
spread and the running basis it implies.
\end{exercise}

\begin{exercise}[$\star\star$]\label{exo:m2:swap-clearing-and-the-ccp-basis:5}
Clients at one house are mostly fixed payers. At which house is the \omterm{def:m2:interest-rate-swaps:par}{par swap
rate} higher, and why?
\end{exercise}

\begin{exercise}[$\star\star$]\label{exo:m2:swap-clearing-and-the-ccp-basis:6}
A fund has EUR~9 billion of uncleared derivatives outstanding and a margin model
that gives EUR~70 million of initial margin against one dealer. Is it subject to
the \omterm{def:m2:swap-clearing-and-the-ccp-basis:umr}{uncleared margin rules}? How much initial margin is exchanged?
\end{exercise}

\begin{exercise}[$\star\star\star$]\label{exo:m2:swap-clearing-and-the-ccp-basis:7}
\emph{Coding.} With \texttt{basis\_bp}, give the thirty-year basis per one-way
position at funding spreads of 25, 50 and 75 basis points. Compare the total
for two houses at 50 with the 3.4 basis points found in 2017.
\end{exercise}

\begin{exercise}[$\star\star\star$]\label{exo:m2:swap-clearing-and-the-ccp-basis:8}
\emph{Find the flaw.} ``The thirty-year swap is 3 basis points higher at one
house than at the other: receive there, pay at the other, and pocket 3 basis
points a year risk-free.'' Correct it.
\end{exercise}

\section{Problem: The One-Way Book}

\begin{problem}[{Weekend problem --- what a dealer charges for a lopsided house}]\label{pb:m2:swap-clearing-and-the-ccp-basis:1}
A dealer has received fixed on USD~2 billion of ten-year swaps from clients at
house A, where clients almost never receive. It hedges by paying fixed on
USD~2 billion at house B. Use the stylised model of
\cref{prop:m2:swap-clearing-and-the-ccp-basis:mva}: 7 basis points a day, five
days, $z = 2.326$, a flat 4\% curve, a funding spread of 50 basis points.

\medskip\noindent\textbf{Part I --- House A.}
\begin{enumerate}
\item Give the position's DV01 at house A.
\item Give its initial margin.
\item Give the variation margin for a 5-basis-point rise in rates, and who pays.
\item Why does the house not care that the dealer is hedged elsewhere?
\item How would the margin change if the dealer's position at A were netted by
receivers from other dealers?
\end{enumerate}

\medskip\noindent\textbf{Part II --- Both houses.}
\begin{enumerate}[resume]
\item Give the initial margin at house B.
\item Give the MVA of each position and of both.
\item Give the running basis that pays for one position, and for both.
\item What would it be for thirty-year swaps?
\item What would initial margin be with a ten-day margin period?
\end{enumerate}

\medskip\noindent\textbf{Part III --- The choice.}
\begin{enumerate}[resume]
\item The dealer could instead hedge at house A, if another dealer would receive
fixed there. What margin would it then post?
\item What is the most the dealer should pay another dealer, in basis points of
rate, to receive fixed from it at A rather than hedge at B?
\item Why might no dealer take that trade?
\item How do the clients end up paying?
\item What would change the imbalance at house A?
\end{enumerate}

\medskip\noindent\textbf{Part IV --- Judgement.}
\begin{enumerate}[resume]
\item Why is margin posted in cash costly at all?
\item Why do real houses' margins differ from this model?
\item How does a rule requiring an account at a given house move the basis?
\item State the \emph{named result}: the basis at which hedging at house B and
netting at house A cost the same, for ten-year swaps.
\item In one sentence: why do identical swaps have two prices?
\end{enumerate}
\end{problem}

\section{Interview questions}

\begin{interviewq}[$\star$ \iqroles{bank, trader}]\label{iq:m2:swap-clearing-and-the-ccp-basis:1}
Why were swaps moved into central clearing after 2008, and what did it change
for a dealer?
\end{interviewq}

\begin{interviewq}[$\star$ \iqroles{bank, developer}]\label{iq:m2:swap-clearing-and-the-ccp-basis:2}
What is the difference between initial and variation margin?
\end{interviewq}

\begin{interviewq}[$\star\star$ \iqroles{trader, bank}]\label{iq:m2:swap-clearing-and-the-ccp-basis:3}
Explain the CME--LCH basis. Is it an arbitrage?
\end{interviewq}

\begin{interviewq}[$\star\star$ \iqroles{researcher, bank}]\label{iq:m2:swap-clearing-and-the-ccp-basis:4}
What is MVA, and how would you estimate it for a new swap?
\end{interviewq}

\begin{interviewq}[$\star\star$ \iqroles{bank}]\label{iq:m2:swap-clearing-and-the-ccp-basis:5}
A client trades a swap bilaterally under a CSA instead of clearing it. What
differs in its costs and risks?
\end{interviewq}

\begin{interviewq}[$\star\star\star$ \iqroles{developer, researcher}]\label{iq:m2:swap-clearing-and-the-ccp-basis:6}
How would you compute the incremental initial margin of a new trade at a
clearing house, quickly enough to quote a client?
\end{interviewq}
